% ============================================
% 第6章：美丽的初稿——AI降低"做出来"的门槛，抬高"知道好不好"的门槛
% 拱桥 · 判断稀缺
% ============================================

\chapter{美丽的初稿}

\vspace{0.5cm}

\begin{center}
{\small\color{muted} 拱桥的秘密，不是抵抗压力，而是利用压力。}
{\small\color{muted} AI把``做出来''的压力降为零——于是``知道好不好''的压力，}
{\small\color{muted} 全部压在了人身上。}
\end{center}

\vspace{1cm}

\section{美丽的初稿}

2026年6月，师宇在知识星球上写了一篇复盘，讲的是一场面向铁路局员工的AI实操培训。他写了一段话：

\begin{shiyu-inline}
``核心概念是'美丽的初稿'：AI不应被理解为替人写完，而是给人一个结构完整、方向正确、可被修改和注入真实经验的初稿。AI写出的是初稿，漂亮、完整、体面，但它还不是你的作品。你要把真实经验、岗位责任、个人判断和行动承诺放进去。否则它只是AI写出的文字，不是你的心得。''

\hfill ——师宇博文，2026年6月12日
\end{shiyu-inline}

\textbf{美丽的初稿。}这四个字，是我认为对AI最准确的定位。

AI给你的东西，不是垃圾。它漂亮、完整、体面。格式规范，逻辑清晰，甚至还附带了参考文献——虽然那些参考文献可能有一半是它编的。它看起来像一个成品。但它不是。

它是\textbf{初稿}。一个结构完整、方向正确、但还没有被注入真实经验的初稿。它需要一个人——一个有真实经验、岗位责任、个人判断和行动承诺的人——把它变成作品。

\section{AI会编造不存在的国家标准}

在同一场培训里，学员提出了一个很具体的问题：

\begin{shiyu-inline}
``AI会编造不存在的国家标准，生成内容会引用已经被替代的旧标准，时效性无法保证。针对细分专业领域，AI生成的内容和实际专业要求不匹配，专业内容准确性较差。''

\hfill ——2026年6月11日，兰州铁路局AI培训，录音记录
\end{shiyu-inline}

AI会编造国家标准。不是因为它``想骗你''——它没有``想''这个功能。是因为它的目标不是``准确''，而是``合理''。它要生成一个看起来最像``国家标准''的东西，而不管那个标准是否真的存在。

\textbf{AI的'错误'不是语法错误。是工程判断错误。}

它不知道一根梁多高才合理。它不知道一个安全系数取多少才够。它不知道一个节点在什么位置才安全。它只是算出了——在它见过的所有文本里——这个数值最像``正确的答案''。

\section{``做出来''的门槛降为零，``知道好不好''的门槛抬向无限}

2026年6月3日，师宇写了一篇博文，标题叫``Skill不是取经路''。里面有段话：

\begin{shiyu-inline}
``AI把'做出来'的门槛降得很低，却把'知道好不好'的门槛抬得更高。Skill可能拉平动作表层，却拉远判断深层。它给普通人专家的外壳，也让真正专家更快迭代。''

\hfill ——师宇博文，2026年6月3日
\end{shiyu-inline}

这是全书最重要的一句话。

\textbf{AI把``做出来''的门槛降得很低。}

以前，你要画一张桥梁效果图，你需要学PS、学渲染、学建模。现在，你输入一句话，AI给你一张图。以前，你要写一份公文，你需要学格式、学模板、学措辞。现在，你输入一个主题，AI给你一份格式规范的草稿。以前，你要搭一座桥的有限元模型，你需要学MIDAS、学ANSYS、学网格划分。现在——你仍然需要。但AI可以帮你生成第一版。

\textbf{但AI把``知道好不好''的门槛抬得更高。}

当所有人都会``做出来''，``做出来''就不再是稀缺能力。稀缺的是``判断''——判断这张图哪里不对，判断这份公文哪里不合适，判断这个模型哪里不安全。而判断，比做出来难得多。

\textbf{做出来只需要一个AI。判断需要一个人。一个被训练过、有经验、有责任感、会在自己错的时候发现自己错的人。}

\begin{insight}
\label{insight:chap6}

\textbf{AI时代最稀缺的能力，不是``写得快''，而是``审得好''。}

不是``画得漂亮''，而是``知道哪里不对''。不是``算得快''，而是``知道这个结果在物理上是否可能''。

生成能力在贬值。判断能力在升值。这就是拱桥的压力转化——AI把生成的压力降为零，把判断的压力全部压在人的身上。而人，恰好就是被这种压力锻造出来的。
\end{insight}

\section{拱桥的隐喻之三}

拱桥的每一块石头都承受着巨大的压力。但正是这种压力，让石头被压紧，让拱桥站住。

\textbf{AI就是那个压力。}

它把``做出来''变得不值钱——于是你被迫去学``判断''。它把``知识传递''变得不值钱——于是你被迫去学``制造困惑''。它把``批改作业''变得不值钱——于是你被迫去学``追问理解''。

这些压力，不是要摧毁你的敌人。它们是建造拱桥的石头。每一块压力，都对应着一块新的能力。就像AI把做出来的门槛降为零，恰好把判断的门槛抬向无限——而判断，正是人最不能被替代的东西。

\begin{shiyu-note}
\label{note:chap6}

\textbf{会生成不稀缺。会审查、会判断、会修改、会提出更高标准才稀缺。}

AI时代的教育，最危险的不是``学生用AI写作业''。最危险的是``学生不会判断AI写的东西对不对''。

作业可以AI写。但判断——判断AI写的东西哪里不对、哪里不够好、哪里需要改——这个能力，AI替不了。而培养这个能力，恰好是教育最应该做的事。

不是禁止AI。是让学生在AI生成的初稿上，练习判断。让他们指出AI哪里错了。让他们修改AI生成的东西。让他们在AI的基础上，做出更好的版本。让他们学会说：``这不是我想要的。我要的是——''

那个``我要的是''，才是教育真正的产出。
\end{shiyu-note}

\vspace{1cm}

\begin{decision-point}
\label{decision:chap6}

\noindent\textbf{这一章，我们看到了AI如何把``做出来''变便宜，把``判断''变贵。}

\vspace{0.3cm}

\noindent 下一节课，做一件事：用AI生成一份你课程里某个知识点的``标准答案''。在课堂上展示给学生。然后问他们：``这是AI给的答案。你们觉得对吗？如果不对，哪里不对？''

\noindent 不要接受``我觉得对''或``我觉得不对''。要求他们\noindent指出\textbf{具体哪里不对}。要求他们\noindent给出\textbf{修改方案}。

\noindent 你会发现，``判断''比``知道''难得多。而那个能指出具体错误的学生，才是真正理解了的人。
\end{decision-point}

\vspace{0.5cm}

\begin{flushright}
{\small\color{muted} 第二部分结束。下一部分：斜拉桥——每一根索都直接传递力，学习回路的设计。}
\end{flushright}